% 07_heuristicas.tex - Capítulo 7: Heurísticas (F3)
\chapter{Heurísticas para A* y Beam Search}
\label{ch:heuristicas}

\section{Requisitos (RF-21, RF-22)}
\label{sec:requisitos-heuristicas}

El enunciado exige \textbf{al menos dos heurísticas} para A*:

\begin{enumerate}
\item \textbf{Heurística 1 (h1):} Distancia Manhattan (obligatoria)
\item \textbf{Heurística 2 (h2):} Propuesta y justificada por el equipo (RF-21)
\end{enumerate}

Para cada heurística se debe justificar (RF-22):
\begin{enumerate}
\item Qué intenta estimar
\item Qué información del estado utiliza
\item Cuál es su costo computacional
\item Si se considera admisible y por qué
\item En qué escenarios orienta bien y en cuáles pierde capacidad de discriminación
\end{enumerate}

Ambas heurísticas se escalan por \code{min\_step\_cost()} (costo transitable mínimo del JSON)
para garantizar admisibilidad aunque el mapa externo traiga costos $< 1$ (enunciado §18.4).

\section{h1: Manhattan Escalada}
\label{sec:h1}

\begin{minted}{python}
def manhattan(problem: PathProblem) -> Heuristic:
    scale = min_step_cost(problem.board)
    goal = problem.goal
    def h(state: GameState) -> float:
        pos = state.position_of(problem.unit_id)
        return (abs(pos.row - goal.row) + abs(pos.col - goal.col)) * scale
    return h
\end{minted}

\begin{table}[H]
\centering
\caption{Justificación RF-22 para h1 (Manhattan)}
\begin{tabular}{lp{8cm}}
\toprule
\textbf{Punto RF-22} & \textbf{Respuesta} \\
\midrule
1. Qué estima & Costo restante si no hay obstáculos y cada paso cuesta el mínimo \\
2. Info del estado & Posición de la unidad, meta, \code{min\_step\_cost} \\
3. Costo computacional & O(1) por evaluación, sin preprocesamiento \\
4. Admisible & Sí. Cada movimiento cambia la distancia Manhattan en $\le 1$ y cuesta $\ge$ min\_cost \\
5. Dónde pierde & Detrás de barreras: apunta recto al objetivo aunque haya muro (caso 3) \\
\bottomrule
\end{tabular}
\end{table}
\label{tab:h1-rf22}


\section{h2: Obstacle-Aware (BFS Inverso desde la Meta)}
\label{sec:h2}

Propuesta del equipo (RF-21): una heurística que \textbf{conoce el layout del mapa} (muros y
base rival) pero ignora costos de terreno y otras unidades.

\begin{minted}{python}
def obstacle_aware(problem: PathProblem) -> Heuristic:
    board = problem.board
    scale = min_step_cost(board)
    rival_base = board.bases[board.unit_sides[problem.unit_id].opponent]
    moves: dict[Position, int] = {problem.goal: 0}
    queue: deque[Position] = deque([problem.goal])
    while queue:
        cell = queue.popleft()
        for nxt in cell.neighbors():
            if nxt in moves or nxt == rival_base or not board.is_passable(nxt):
                continue
            moves[nxt] = moves[cell] + 1
            queue.append(nxt)
    def h(state: GameState) -> float:
        steps = moves.get(state.position_of(problem.unit_id))
        return math.inf if steps is None else steps * scale
    return h
\end{minted}

Construcción: \textbf{BFS inverso desde la meta} (backward BFS) esquivando muros y base rival.
Complejidad: O(R$\times$C) una vez al construir; O(1) por evaluación (lookup en dict).

\begin{table}[H]
\centering
\caption{Justificación RF-22 para h2 (Obstacle-aware)}
\begin{tabular}{lp{8cm}}
\toprule
\textbf{Punto RF-22} & \textbf{Respuesta} \\
\midrule
1. Qué estima & Costo restante con layout real (muros + base rival), paso a costo mínimo \\
2. Info del estado & Posición de la unidad + layout completo del mapa (muros, base rival) \\
3. Costo computacional & O(R$\times$C) una vez (BFS inverso), O(1) por evaluación (lookup) \\
4. Admisible & Sí. Es el costo exacto de un problema relajado (movimientos gratis excepto min\_cost) \\
5. Dónde pierde & No ve costos de terreno: si el atajo cruza pantano, h2 sigue prefiriéndolo (caso 4) \\
\bottomrule
\end{tabular}
\end{table}
\label{tab:h2-rf22}


\section{Propiedades: h2 Domina a h1}
\label{sec:h2-domina}

\begin{itemize}
\item h2 $\ge$ h1 \textbf{siempre} (h2 cuenta movimientos reales esquivando obstáculos; h1 asume camino recto)
\item Ambas son \textbf{admisibles} y \textbf{consistentes} para movimientos de 4 direcciones
\item h2 tiene O(R$\times$C) preprocesamiento una vez; h1 es O(1) sin preprocesamiento
\item En A*, h2 expande \textbf{menos estados} que h1 (casos 3 y 4 lo demuestran)
\item Celdas inalcanzables $\to$ h2 = $\infty$ (admisible: no hay camino, costo real = $\infty$)
\end{itemize}

\section{Medición de Efecto (Casos 3 y 4)}
\label{sec:medicion-heuristicas}

\begin{table}[H]
\centering
\caption{Expandidos en Caso 3 (muro) y Caso 4 (pantano vs. camino abierto)}
\begin{tabular}{lccc}
\toprule
Algoritmo & Caso 3 (muro) & Caso 4 (pantano vs. abierto) & Nota \\
\midrule
UCS & 295 & 332 & baseline \\
A* h1 (Manhattan) & 166 & 186 & h1 ignora muro/pantano \\
A* h2 (Obstacle-aware) & 84 & 115 & h2 ve muro, no ve pantano \\
\bottomrule
\end{tabular}
\end{table}
\label{tab:heuristicas-expandidos}


En caso 3 (muro vertical), h2 reduce expandidos a la \textbf{mitad} de h1 (84 vs 166).
En caso 4 (atajo por pantano), h2 ayuda pero no ve el costo del terreno.

\section{Origen de h2}
\label{sec:origen-h2}

La heurística h2 fue \textbf{propuesta por Claude Code (Sonnet 5.5)} durante el desarrollo;
el equipo revisó la admisibilidad, la consistencia y la dominancia sobre h1, y la aceptó
tras verificar que cumple los 5 puntos de RF-22. (La documentación inicial la llamaba
``propuesta del equipo''; se corrigió en \texttt{decisions.md} 2026-10-17).